% =====================================================================
%  30.3  The Matrix-Tree theorem.
%  Exports: sec:c32-matrix  def:c32-lap  thm:c32-matrix  lem:c32-cor
%  lst:c32-lap  lst:c32-moddet  fig:c32-lap  fig:c32-c5
%  WORKED examples only -- no solved problem here.  The Matrix-Tree
%  theorem gives spanning-tree count as any cofactor of the Laplacian;
%  a modular determinant is the whole algorithm for large counts.
%  Sanity checks against Cayley (K_n -> n^{n-2}), cycle C_n -> n, and
%  complete bipartite K_{a,b} -> a^{b-1} b^{a-1}.
% =====================================================================
\section{The Matrix-Tree Theorem}\label{sec:c32-matrix}

\begin{previouslybox}[Pr\"ufer handled the complete graph; now the graph is given]
\(\secref{c32-prufer}\) and \(\secref{c32-cayley}\) enumerate \emph{all}
labelled trees on \([n]\) or those satisfying a degree filter.  The input of
this section is different: an \emph{explicit} graph \(G=(V,E)\), and the
question is how many spanning trees \(G\) has.  When \(G=K_n\) the answer is
Cayley's \(n^{n-2}\); when \(G\) is sparser you cannot reduce to a Pr\"ufer
code filter -- you need a different tool.  The Matrix-Tree theorem
(\(\thmref{c32-matrix}\)) is that tool: a single determinant of a matrix you
can build in \(\bigO{n+m}\) gives the answer.
\end{previouslybox}

\begin{topicintro}
Attach to \(G\) its \term{Laplacian} \(L = D - A\) (\(\defref{c32-lap}\)), the
\(n\times n\) matrix with \(L_{uu}=\deg(u)\) on the diagonal and \(L_{uv} =
-\)(number of edges \(u\!-\!v\)) off.  Delete row and column \(i\) for any
fixed \(i\); the resulting \((n-1)\times(n-1)\) cofactor \(L^{(i)}\) has a
determinant equal to the number of spanning trees of \(G\)
(\(\thmref{c32-matrix}\)), independent of which \(i\) you chose.  For
unweighted graphs the count is an integer and can be huge; work modulo a
prime \(p\) and evaluate the determinant by Gaussian elimination in
\(\bigO{n^3}\) -- that is \(\coderef{c32-moddet}\).

The section is \emph{worked examples only}: no CF problem is solved here,
because the pure Matrix-Tree template almost always appears bundled with a
graph-construction twist (a problem you meet in Ch.~29 or Ch.~31).  What is
here is the mathematical engine itself and three sanity checks: cycle \(C_n\)
gives \(n\); complete bipartite \(K_{a,b}\) gives \(a^{b-1}b^{a-1}\);
complete graph \(K_n\) gives \(n^{n-2}\) recovering Cayley.

\textbf{Prerequisites.} Modular inverse (Ch.~6 background) via Fermat's little
theorem for a prime modulus; standard Gaussian elimination on a matrix.
Nothing else is assumed.

\textbf{After this section you can}
\begin{itemize}[nosep]
  \item build \(L(G)\) in \(\bigO{n+m}\) directly from an adjacency list; \hfill\lv{3}
  \item state the theorem and its dependence on the all-ones kernel of \(L\); \hfill\lv{4}
  \item evaluate the modular determinant of \(L^{(i)}\) via \(\coderef{c32-moddet}\); \hfill\lv{4}
  \item sanity-check the answer against Cayley and known families. \hfill\lv{4}
\end{itemize}
\end{topicintro}

% ---------------------------------------------------------------
\subsection{The Laplacian and the cofactor}
% ---------------------------------------------------------------

\begin{definition}{Laplacian matrix}{c32-lap}
Let \(G=(V,E)\) be a finite multigraph on \(V=[n]\) (parallel edges allowed;
self-loops ignored in \(L\)).  Its \term{Laplacian} is the \(n\times n\)
integer matrix
\[
  L(G)_{uv} \;=\; \begin{cases}
    \deg_G(u), & u=v,\\
    -(\text{number of edges }u\!-\!v), & u\ne v.
  \end{cases}
\]
Equivalently, \(L = D - A\) with \(D\) the degree diagonal and \(A\) the
adjacency matrix.
\end{definition}

\begin{lemma}{Kernel of the Laplacian}{c32-cor}
For every \(G\), \(L(G)\cdot \mathbf{1} = 0\) (each row of \(L\) sums to
zero).  If \(G\) is connected, the null-space of \(L(G)\) is exactly the
span of \(\mathbf{1}\), so \(\mathrm{rank}\,L(G) = n-1\); deleting any one
row and the corresponding column gives a nonsingular matrix.
\end{lemma}

\begin{proofbox}
Row \(u\) of \(L\) has diagonal \(\deg(u)\) and off-diagonals \(-1\) per edge
incident to \(u\); summing across the row contributes \(\deg(u) - \deg(u) =
0\).  If \(x\in\ker L\) then \(x^\top L x = \sum_{(u,v)\in E}(x_u-x_v)^2\);
this vanishes iff \(x_u=x_v\) on every edge, i.e.\ iff \(x\) is constant on
each connected component.  Connectivity forces \(x\propto\mathbf{1}\).
\(\Box\)
\end{proofbox}

\begin{theorem}{Matrix-Tree theorem (Kirchhoff)}{c32-matrix}
For a connected graph \(G\) on \([n]\), the number of its spanning trees
\(\tau(G)\) equals \(\det L^{(i)}(G)\) for any \(i\in[n]\), where \(L^{(i)}\)
is the \((n-1)\times(n-1)\) matrix obtained by deleting row and column \(i\)
from \(L(G)\).  For a disconnected \(G\), every \((n-1)\)-cofactor has
determinant zero (there are no spanning trees).
\end{theorem}

The theorem has four standard proofs -- via the Cauchy-Binet identity
applied to a reduced incidence matrix, via a minors-of-a-minor argument on
\(L\), via an electrical-network effective-resistance identity, and via a
combinatorial species bijection.  We give the Cauchy-Binet one, since it
builds directly on \(\lemref{c32-cor}\).

\begin{proofbox}
Pick an orientation of \(G\)'s edges, and let \(B\) be its \((n-1)\times m\)
reduced incidence matrix (row \(u\) is the incidence vector of edge \(u\)
minus the reference row \(i\)).  Then \(L^{(i)} = B\,B^\top\) (a standard
identity you can verify entrywise: the diagonal counts incident edges, the
off-diagonal subtracts the shared-edge indicator).  By Cauchy-Binet, \(\det
(BB^\top) = \sum_{S\subseteq[m], |S|=n-1} (\det B_S)^2\), where \(B_S\) is the
square submatrix on edge set \(S\).  A classical fact about \(B\) is that
\(\det B_S\in\{+1,-1\}\) iff \(S\) is the edge set of a spanning tree of
\(G\), and \(\det B_S = 0\) otherwise.  So each spanning tree contributes
exactly \(1\) to the sum, and non-tree subsets contribute \(0\).  Hence
\(\det L^{(i)} = \tau(G)\). \(\Box\)
\end{proofbox}

\begin{lstlisting}[caption={Build the Laplacian of a graph on n vertices with m edges, then delete row and column 0 (equivalent to any other by the theorem). Modular Gaussian elimination evaluates the determinant. ASCII only, identifiers without underscores}, label={lst:c32-lap}]
#include <bits/stdc++.h>
using namespace std;
typedef long long ll;
const ll MOD=1000000007LL;
ll modPow(ll b, ll e, ll m){
  ll r=1; b%=m;
  while(e>0){ if(e&1) r=(r*b)%m; b=(b*b)%m; e>>=1; }
  return r;
}
ll modInv(ll a){ return modPow(a, MOD-2, MOD); }
int main(){
  int n,m; if(scanf("%d %d",&n,&m)!=1) return 0;
  vector<vector<ll> > L(n, vector<ll>(n,0));
  for(int i=0;i<m;i++){
    int u,v; scanf("%d %d",&u,&v); u--; v--;
    L[u][u]++; L[v][v]++;
    L[u][v] = (L[u][v]-1+MOD)%MOD;
    L[v][u] = (L[v][u]-1+MOD)%MOD;
  }
  int sz=n-1;
  vector<vector<ll> > M(sz, vector<ll>(sz,0));
  for(int i=0;i<sz;i++) for(int j=0;j<sz;j++) M[i][j]=L[i+1][j+1];
  ll det=1;
  for(int col=0; col<sz; col++){
    int piv=-1;
    for(int r=col; r<sz; r++) if(M[r][col]%MOD!=0){ piv=r; break; }
    if(piv<0){ det=0; break; }
    if(piv!=col){ swap(M[piv], M[col]); det=(MOD-det)%MOD; }
    det=(det*M[col][col])%MOD;
    ll inv=modInv(M[col][col]);
    for(int r=col+1; r<sz; r++){
      ll factor = (M[r][col]*inv)%MOD;
      for(int c=col; c<sz; c++)
        M[r][c] = (M[r][c] - factor*M[col][c]%MOD + MOD)%MOD;
    }
  }
  printf("%lld\n", det);
  return 0;
}
\end{lstlisting}

\begin{lstlisting}[caption={Standalone modular determinant of an \(s\times s\) matrix under a prime \(p\), used repeatedly when the same Laplacian is queried for different minors. Gaussian elimination with row swap and multiplicative pivots; \(\bigO{s^3}\)}, label={lst:c32-moddet}]
ll modDet(vector<vector<ll> > M, ll p){
  int s=(int)M.size();
  ll det=1;
  for(int col=0; col<s; col++){
    int piv=-1;
    for(int r=col; r<s; r++) if(M[r][col]%p!=0){ piv=r; break; }
    if(piv<0) return 0;
    if(piv!=col){ swap(M[piv], M[col]); det=(p-det)%p; }
    det=(det*M[col][col])%p;
    ll inv=modPow(M[col][col], p-2, p);
    for(int r=col+1; r<s; r++){
      ll factor=(M[r][col]*inv)%p;
      for(int c=col; c<s; c++)
        M[r][c]=(M[r][c]-factor*M[col][c]%p+p)%p;
    }
  }
  return det;
}
\end{lstlisting}

% ---------------------------------------------------------------
\subsection{Worked examples: three sanity checks}
% ---------------------------------------------------------------

\subsection*{Worked 1: The cycle \(C_5\)}

\(C_5\) on labels \(1,2,3,4,5\) with edges \((1,2),(2,3),(3,4),(4,5),(5,1)\).
Its Laplacian is tridiagonal on the ring, degree \(2\) on the diagonal and
\(-1\) on the cyclic off-diagonals.  Delete row and column \(5\); the
remaining \(4\times 4\) matrix is the path Laplacian
\[
  L^{(5)}(C_5) = \begin{pmatrix} 2 &-1 & 0 & 0\\ -1 & 2 &-1 & 0\\ 0&-1&2&-1\\ 0&0&-1&2\end{pmatrix}.
\]
Its determinant is the Fibonacci-like sequence \(d_1=2,\, d_2=3,\, d_3=4,\,
d_4=5\) (each cofactor of a path Laplacian on \(k\) nodes with fixed
boundaries is \(k+1\)); the answer is \(5\).  And indeed \(C_5\) has
\emph{exactly} \(5\) spanning trees -- remove any one of the five edges.
The formula matches the combinatorics.

\subsection*{Worked 2: The complete bipartite graph \(K_{2,3}\)}

Partition \(A=\{1,2\}\), \(B=\{3,4,5\}\), every \(A\)-vertex joined to every
\(B\)-vertex; no intra-part edges.  Degrees: \(\deg(1)=\deg(2)=3,\)
\(\deg(3)=\deg(4)=\deg(5)=2\).  Delete row and column \(5\); the resulting
\(4\times 4\) cofactor on \(\{1,2,3,4\}\) is
\[
  L^{(5)} = \begin{pmatrix}3&0&-1&-1\\0&3&-1&-1\\-1&-1&2&0\\-1&-1&0&2\end{pmatrix}.
\]
Reduce: rows 3 and 4 differ by \((0,0,2,-2)\), giving a pivot at \((3,3)=2\);
rows 1 and 2 by \((3,0,-1,-1)\) versus \((0,3,-1,-1)\), independent.  A
short elimination produces \(\det = 12\).  The closed formula
\(a^{b-1}b^{a-1}=2^{2}\cdot 3^{1}=12\) agrees.

\subsection*{Worked 3: The complete graph \(K_4\)}

\(L(K_4)=4I - (J-I) = 4I - J + I = 5I - J\) with \(J\) the all-ones matrix.
Delete row/col 4 to get \(L^{(4)}=5I_3 - J_3 = \mathrm{diag}(5)-\mathrm{off\text{-}diag}(1)\).
Its eigenvalues are \(5-3=2\) (once, on \(\mathbf{1}\)) and \(5\) (twice, on
the orthogonal complement), so the determinant is \(2\cdot 5\cdot 5=50\)??  No
-- re-check.  \(L(K_4)\) has diagonal \(3\) (degree \(3\) in \(K_4\)), not
\(5\).  Correct matrix \(L^{(4)}\) on the top-left \(3\times 3\) block has
diagonal \(3\) and off-diagonal \(-1\), so \(L^{(4)}=4I_3 - J_3\).
Eigenvalues on the same split: \(4-3=1\) on \(\mathbf{1}\) and \(4\) twice on
the orthogonal complement, so \(\det = 1\cdot 4 \cdot 4 = 16\).  And indeed
Cayley \(n^{n-2}=4^2=16\).  The formula matches \(\thmref{c32-cayley}\).

\begin{figure}[H]
\centering
\begin{tikzpicture}[diagram, scale=0.9, transform shape,
   tnode/.style={circle, draw=primaryblue!70, fill=primaryblue!8, line width=0.7pt,
                 minimum size=6mm, font=\scriptsize\bfseries, inner sep=1pt}]
  \node[tnode] (a) at (0,1) {1};
  \node[tnode] (b) at (1.2,1) {2};
  \node[tnode] (c) at (0,0) {3};
  \node[tnode] (d) at (1.2,0) {4};
  \node[tnode] (e) at (0.6,-1) {5};
  \draw (a)--(c)--(e)--(d)--(b)--(a);
  \draw (a)--(b); \draw (c)--(d);
  \node[annot, font=\scriptsize, anchor=north] at (0.6,-2.4)
    {\(C_5\) on five labels: delete any row/col pair and the \(4\times 4\) cofactor has determinant \(5\)};
\end{tikzpicture}
\caption{The five-cycle \(C_5\); \(\tau=5\) because a spanning tree is
exactly the graph minus one of the five edges.  Worked check of
\(\thmref{c32-matrix}\) against the combinatorial count.}
\label{fig:c32-c5}
\end{figure}

\begin{insightbox}[What Matrix-Tree buys you that Pr\"ufer does not]
Pr\"ufer's bijection is a closed form on \(\mathcal{T}_n\); it is unbeatable
when the answer is a sequence enumeration.  Matrix-Tree is the \emph{generic
engine} -- hand it any explicit graph and it returns the count.  The two
tools are not rivals; the problems that mix them (given a graph, count its
spanning trees \emph{with a degree constraint}) are exactly the CF 2400+ tier
where you must decide which tool handles which part of the constraint.
Section \(\secref{c32-boundary}\) makes the split precise, and the Summit
\(\probref{c32-summit}\) uses the completion formula from
\(\secref{c32-cayley}\) rather than the determinant because the graph there
is a \emph{given tree}, not an arbitrary one.
\end{insightbox}

\begin{notebox}[Modular pitfalls]
\textbf{Modular pitfalls.}  The determinant algorithm assumes the modulus is
prime (or at least that every pivot is invertible).  Use \(p=10^9+7\) or
\(p=998244353\); avoid composite moduli unless you switch to Bareiss' exact
division (slower but works over \(\mathbb{Z}\)).  Row swaps flip the sign,
handled by \((\mathrm{MOD}-\det)\bmod\mathrm{MOD}\).  Zero pivots on a
nonzero determinant are handled by picking a row below to swap in.
\end{notebox}
